\chapter{模论}

本章进入环论之后最重要的对象之一：模。在线性代数中，向量空间是定义在域上的“线性对象”；模则是把标量域推广为一般环之后得到的自然概念。它既保留了线性代数的许多思想，又允许出现更丰富、更精细的现象。对模的理解，是学习交换代数、代数数论、表示论以及模形式理论所必需的基础。

除特别说明外，本章中的环都默认含有单位元，且环同态保持单位元；模默认是\emph{左模}。若 $R$ 为交换环，则左右模没有本质区别，我们便直接称之为 $R$-模。

\section{模的定义与例子}

\subsection{模与模同态}

\begin{definition}
设 $R$ 为环。一个\emph{$R$-模}是一个阿贝尔群 $(M,+)$，并赋予一个数量乘法
\[
R\times M\longrightarrow M,\qquad (r,m)\longmapsto rm,
\]
满足对任意 $r,s\in R$ 与 $m,n\in M$，有
\begin{align*}
(r+s)m&=rm+sm,\\
r(m+n)&=rm+rn,\\
(rs)m&=r(sm),\\
1_Rm&=m.
\end{align*}
\end{definition}

这些公理说明：模既是一个加法阿贝尔群，又允许环中的元素作为“标量”对其作用。若把 $R$ 换成域 $F$，则这正是向量空间的定义。

\begin{remark}
有些文献在讨论不含单位元的环或非幺模时，会去掉条件 $1_Rm=m$。本书始终采用含单位元的情形，因为这才是大多数代数与数论应用中的自然框架。
\end{remark}

\begin{definition}
设 $M,N$ 为 $R$-模。映射 $f\colon M\to N$ 称为\emph{$R$-模同态}，若对任意 $m,m'\in M$ 与 $r\in R$，都有
\[
f(m+m')=f(m)+f(m'),\qquad f(rm)=rf(m).
\]
若 $f$ 双射，则称 $f$ 为\emph{$R$-模同构}，记 $M\cong N$。
\end{definition}

模同态的核与像分别定义为
\[
\Ker f=\{m\in M:f(m)=0\},\qquad \Ima f=f(M).
\]
与群同态、线性映射的情形一样，核测量“退化”，像描述“达到的部分”。

\begin{proposition}
设 $f\colon M\to N$ 为 $R$-模同态，则 $\Ker f$ 与 $\Ima f$ 都分别是 $M$ 与 $N$ 的子模。
\end{proposition}

\begin{proof}
我们先证 $\Ker f$。因为 $f(0)=f(1_R\cdot 0)=1_Rf(0)$，故 $f(0)=0$，从而 $0\in \Ker f$。若 $x,y\in \Ker f$，则
\[
f(x-y)=f(x)-f(y)=0-0=0,
\]
故 $x-y\in \Ker f$。再对任意 $r\in R$ 与 $x\in \Ker f$，有
\[
f(rx)=rf(x)=r\cdot 0=0,
\]
故 $rx\in \Ker f$。于是 $\Ker f$ 是子模。

再证 $\Ima f$。显然 $0=f(0)\in \Ima f$。若 $u=f(x),v=f(y)\in \Ima f$，则
\[
u-v=f(x)-f(y)=f(x-y)\in \Ima f.
\]
对任意 $r\in R$，有
\[
ru=rf(x)=f(rx)\in \Ima f.
\]
故 $\Ima f$ 也是子模。
\end{proof}

\subsection{基本例子}

\begin{example}[阿贝尔群就是 $\Z$-模]
任一阿贝尔群 $A$ 都自然成为 $\Z$-模：对 $n\in \Z$ 与 $a\in A$，定义
\[
na=\begin{cases}
\underbrace{a+\cdots+a}_{n\text{ 次}},&n>0,\\
0,&n=0,\\
-( (-n)a),&n<0.
\end{cases}
\]
模公理由整数乘法与阿贝尔群运算的基本性质立刻得到。反过来，任意 $\Z$-模的加法结构就是一个阿贝尔群。因此“阿贝尔群范畴”等同于“$\Z$-模范畴”。
\end{example}

\begin{example}[向量空间就是域上的模]
设 $F$ 为域。$F$-模正是 $F$-向量空间，因为模公理与线性代数中向量空间的公理完全一致。因而模论可以看作线性代数在一般环上的推广。
\end{example}

\begin{example}[理想是自然的模]
设 $R$ 为环，$I\subseteq R$ 为左理想。则 $I$ 在加法下是阿贝尔群，并且对任意 $r\in R$ 与 $x\in I$，有 $rx\in I$，故 $I$ 是一个左 $R$-模。若 $R$ 交换，则任意理想都是 $R$-模。
\end{example}

\begin{example}[环自身作为模]
环 $R$ 在自身加法下是阿贝尔群，并以环乘法作为数量乘法，因此 $R$ 是一个左 $R$-模。若 $R$ 交换，则也可看成右 $R$-模。这个最基本的例子常常把环论问题翻译为模论问题，例如“理想”就是“$R$ 的子模”。
\end{example}

\begin{example}[直积与直和]
若 $M_i$ $(i\in I)$ 是一族 $R$-模，则其直积
\[
\prod_{i\in I} M_i
\]
按分量加法与分量数量乘法成为 $R$-模。若只允许有限多个分量非零，则得到直和
\[
\bigoplus_{i\in I} M_i,
\]
它也是 $R$-模。当 $I$ 有限时，直积与直和一致。
\end{example}

\begin{example}[商环诱导的模]
设 $I$ 为交换环 $R$ 的理想，则商环 $R/I$ 作为阿贝尔群，可以通过
\[
r\cdot (a+I)=ra+I
\]
定义成一个 $R$-模。此定义良定，因为若 $a-a'\in I$，则 $r(a-a')\in I$。这个例子在以后构造商模和张量积时都会反复出现。
\end{example}

\subsection{$R[x]$-模与线性算子}

下面的例子把模论与线性代数联系起来。

\begin{proposition}
设 $F$ 为域，$V$ 为 $F$-向量空间。则在 $V$ 上给定一个 $F[x]$-模结构，与在 $V$ 上给定一个 $F$-线性算子 $T\in\End_F(V)$ 是等价的。
\end{proposition}

\begin{proof}
先设 $V$ 已有 $F[x]$-模结构。因为 $x\in F[x]$，可定义映射
\[
T\colon V\to V,\qquad T(v)=xv.
\]
我们证明 $T$ 是 $F$-线性的。对 $v,w\in V$，有
\[
T(v+w)=x(v+w)=xv+xw=T(v)+T(w).
\]
对任意 $a\in F$，注意 $a$ 可视为常数多项式，故
\[
T(av)=x(av)=(xa)v=(ax)v=a(xv)=aT(v),
\]
这里用到了 $F[x]$ 中常数多项式与 $x$ 可交换。故 $T\in \End_F(V)$。

反过来，设给定 $F$-线性算子 $T\in\End_F(V)$。对任意多项式
\[
f(x)=a_0+a_1x+\cdots+a_nx^n\in F[x],
\]
定义其在 $V$ 上的作用为
\[
f(x)\cdot v\coloneqq a_0v+a_1T(v)+\cdots+a_nT^n(v).
\]
由于 $T$ 线性，上式对 $v$ 是良定义的。我们验证模公理。设 $f,g\in F[x]$ 与 $v,w\in V$。

其一，$(f+g)\cdot v=f\cdot v+g\cdot v$ 由定义立即成立；其二，
\[
f\cdot (v+w)=f\cdot v+f\cdot w
\]
也由 $T^i$ 的线性得到。关键是乘法相容性。写
\[
f(x)=\sum_{i=0}^m a_ix^i,\qquad g(x)=\sum_{j=0}^n b_jx^j.
\]
则
\begin{align*}
(fg)\cdot v
&=\sum_{k=0}^{m+n}\Bigl(\sum_{i+j=k}a_ib_j\Bigr)T^k(v)\\
&=\sum_{i=0}^m\sum_{j=0}^n a_ib_jT^{i+j}(v)\\
&=\sum_{i=0}^m a_iT^i\Bigl(\sum_{j=0}^n b_jT^j(v)\Bigr)\\
&=f\cdot (g\cdot v).
\end{align*}
最后，常数多项式 $1$ 的作用显然是恒等映射。因此得到一个 $F[x]$-模结构。

上述两个过程互为逆过程：从模结构得到的算子是 $T(v)=xv$，再由该算子恢复多项式作用，就得到原模结构；反之亦然。
\end{proof}

\begin{remark}
这个命题说明：研究单个线性算子 $T$ 的问题，可以改写为研究一个 $F[x]$-模。在线性代数中，最小多项式、特征多项式、Jordan 标准形与有理标准形，都可以用模论语言统一表达。
\end{remark}

\section{子模与商模}

\subsection{子模}

\begin{definition}
设 $M$ 为 $R$-模。子集 $N\subseteq M$ 称为 $M$ 的一个\emph{子模}，若
\begin{enumerate}[label=(\arabic*)]
  \item $0\in N$；
  \item 对任意 $x,y\in N$，有 $x-y\in N$；
  \item 对任意 $r\in R$ 与 $x\in N$，有 $rx\in N$。
\end{enumerate}
记作 $N\le M$。
\end{definition}

因此，子模就是既对加法封闭、又对环作用封闭的子集。它同时推广了向量空间中的子空间与群论中的正规子群。

\begin{example}
$\Z$-模 $\Z$ 的子模恰好是 $n\Z$ $(n\ge 0)$。证明如下：设 $N\le \Z$。若 $N=\{0\}$ 则为 $0\Z$。若 $N\ne\{0\}$，则 $N\cap \N$ 非空，由良序原理取其中最小正元 $d$。对任意 $n\in N$，由带余除法有 $n=qd+r$，其中 $0\le r<d$。由于 $N$ 是子模，$r=n-qd\in N$；由 $d$ 的最小性知 $r=0$，故 $n\in d\Z$。反过来 $d\in N$ 蕴含 $qd\in N$，因此 $N=d\Z$。
\end{example}

\begin{example}
若 $V$ 是域 $F$ 上向量空间，则 $V$ 的子模就是它的线性子空间。
\end{example}

\begin{example}
若 $R$ 为交换环，则 $R$ 作为 $R$-模的子模就是它的理想。这一观察使得许多环论陈述都可视为模论陈述的特例。
\end{example}

\subsection{商模}

\begin{definition}
设 $N\le M$。定义等价关系
\[
m\sim m'\iff m-m'\in N.
\]
等价类记为 $m+N$。全体等价类记为 $M/N$，称为\emph{商模}；其加法与标量乘法定义为
\[
(m+N)+(m'+N)=(m+m')+N,
\]
\[
r(m+N)=rm+N.
\]
\end{definition}

\begin{proposition}
上述运算良定义，并使 $M/N$ 成为一个 $R$-模。
\end{proposition}

\begin{proof}
先证良定义。若 $m+N=m_1+N$ 且 $m'+N=m_1'+N$，则 $m-m_1\in N$ 与 $m'-m_1'\in N$。因为 $N$ 是子模，
\[
(m+m')-(m_1+m_1')=(m-m_1)+(m'-m_1')\in N,
\]
故 $(m+m')+N=(m_1+m_1')+N$。同理，对任意 $r\in R$，
\[
rm-rm_1=r(m-m_1)\in N,
\]
故 $rm+N=rm_1+N$。于是运算良定义。

再验证模公理。由于商群 $M/N$ 在加法下已是阿贝尔群，只需检查数量乘法的四条公理；它们都直接来自 $M$ 中对应公理。例如
\[
(r+s)(m+N)=((r+s)m)+N=(rm+sm)+N=(rm+N)+(sm+N),
\]
其余同理。
\end{proof}

\begin{example}
设 $M=\Z$，$N=n\Z$。则商模 $\Z/n\Z$ 的元素为余数类 $\overline{0},\overline{1},\dots,\overline{n-1}$，它是最基本的有限 $\Z$-模。
\end{example}

\begin{definition}
自然映射
\[
\pi\colon M\to M/N,\qquad \pi(m)=m+N,
\]
称为从 $M$ 到 $M/N$ 的\emph{商映射}。
\end{definition}

\begin{proposition}
自然映射 $\pi\colon M\to M/N$ 是满的 $R$-模同态，且 $\Ker \pi=N$。
\end{proposition}

\begin{proof}
对任意 $m,m'\in M$ 与 $r\in R$，有
\[
\pi(m+m')=(m+m')+N=(m+N)+(m'+N)=\pi(m)+\pi(m'),
\]
以及
\[
\pi(rm)=rm+N=r(m+N)=r\pi(m),
\]
故 $\pi$ 为模同态。其满射性显然，因为任意商类都形如 $m+N$。又
\[
\pi(m)=0+N \iff m+N=N \iff m\in N,
\]
故 $\Ker \pi=N$。
\end{proof}

\subsection{同构定理}

\begin{theorem}[第一同构定理]
设 $f\colon M\to N$ 为 $R$-模同态，则
\[
M/\Ker f\cong \Ima f.
\]
更具体地，映射
\[
\overline f\colon M/\Ker f\to \Ima f,\qquad \overline f(m+\Ker f)=f(m),
\]
是一个良定义的 $R$-模同构。
\end{theorem}

\begin{proof}
先证良定义。若 $m+\Ker f=m'+\Ker f$，则 $m-m'\in \Ker f$，故
\[
f(m)-f(m')=f(m-m')=0,
\]
从而 $f(m)=f(m')$，故 $\overline f$ 良定义。

再证其为模同态。对任意 $m,m'\in M$ 与 $r\in R$，有
\begin{align*}
\overline f((m+\Ker f)+(m'+\Ker f))
&=\overline f(m+m'+\Ker f)=f(m+m')\\
&=f(m)+f(m')=\overline f(m+\Ker f)+\overline f(m'+\Ker f),
\end{align*}
并且
\[
\overline f(r(m+\Ker f))=\overline f(rm+\Ker f)=f(rm)=rf(m)=r\,\overline f(m+\Ker f).
\]
因此 $\overline f$ 为模同态。

其像等于 $\Ima f$，故满射。若 $\overline f(m+\Ker f)=0$，则 $f(m)=0$，即 $m\in \Ker f$，从而 $m+\Ker f=0+\Ker f$。故 $\Ker \overline f=0$，即 $\overline f$ 单射。于是 $\overline f$ 是同构。
\end{proof}

\begin{theorem}[第二同构定理]
设 $A,B\le M$，则 $A+B\le M$，$A\cap B\le A$，并且
\[
(A+B)/B\cong A/(A\cap B).
\]
\end{theorem}

\begin{proof}
先证 $A+B$ 为子模。显然 $0=0+0\in A+B$。若 $a_1+b_1,a_2+b_2\in A+B$，则
\[
(a_1+b_1)-(a_2+b_2)=(a_1-a_2)+(b_1-b_2)\in A+B.
\]
对任意 $r\in R$，有
\[
r(a+b)=ra+rb\in A+B.
\]
故 $A+B$ 为子模。$A\cap B$ 是子模也显然。

定义
\[
\varphi\colon A\to (A+B)/B,\qquad \varphi(a)=a+B.
\]
因为 $a\in A\subseteq A+B$，该映射有意义。并且 $\varphi$ 是模同态。其像为全部 $(A+B)/B$：任取 $(a+b)+B\in (A+B)/B$，有
\[
(a+b)+B=a+B=\varphi(a).
\]
故 $\varphi$ 满射。其核为
\[
\Ker\varphi=\{a\in A:a+B=B\}=\{a\in A:a\in B\}=A\cap B.
\]
由第一同构定理得
\[
A/(A\cap B)\cong (A+B)/B.
\]
\end{proof}

\begin{theorem}[第三同构定理]
设 $L\le N\le M$，则 $N/L\le M/L$，且
\[
(M/L)/(N/L)\cong M/N.
\]
\end{theorem}

\begin{proof}
先证 $N/L$ 是 $M/L$ 的子模。因 $N$ 为子模，$0+L\in N/L$；若 $n_1+L,n_2+L\in N/L$，则
\[
(n_1+L)-(n_2+L)=(n_1-n_2)+L\in N/L;
\]
对任意 $r\in R$，
\[
r(n+L)=rn+L\in N/L.
\]
故 $N/L\le M/L$。

定义
\[
\psi\colon M/L\to M/N,\qquad \psi(m+L)=m+N.
\]
若 $m+L=m'+L$，则 $m-m'\in L\subseteq N$，故 $m+N=m'+N$，因此 $\psi$ 良定义。它显然是满的模同态。其核为
\[
\Ker\psi=\{m+L:m+N=N\}=\{m+L:m\in N\}=N/L.
\]
由第一同构定理得
\[
(M/L)/(N/L)\cong M/N.
\]
\end{proof}

\begin{theorem}[对应定理]
设 $N\le M$。则包含 $N$ 的 $M$ 的子模与商模 $M/N$ 的子模之间存在一一对应：
\[
L\longmapsto L/N.
\]
其逆对应为
\[
P\longmapsto \pi^{-1}(P),
\]
其中 $\pi\colon M\to M/N$ 为自然映射。并且在该对应下，包含关系、和与交都被保持：
\[
(L_1+L_2)/N=(L_1/N)+(L_2/N),\qquad (L_1\cap L_2)/N=(L_1/N)\cap (L_2/N).
\]
\end{theorem}

\begin{proof}
若 $L$ 为包含 $N$ 的子模，则 $L/N$ 是 $M/N$ 的子模，这是第三同构定理证明中的同样验证。反过来，若 $P\le M/N$，则 $\pi^{-1}(P)$ 是 $M$ 的子模，且显然包含 $N=\Ker\pi$。

我们证明二者互逆。首先，对包含 $N$ 的子模 $L$，有
\[
\pi^{-1}(L/N)=L,
\]
因为 $m\in \pi^{-1}(L/N)$ 当且仅当 $m+N\in L/N$，这等价于存在 $\ell\in L$ 使 $m+N=\ell+N$，即 $m-\ell\in N\subseteq L$，从而 $m\in L$；反向包含显然。

其次，对 $P\le M/N$，有
\[
\pi(\pi^{-1}(P))=P,
\]
因为 $\pi$ 满射。于是得到一一对应。

关于运算保持性，取包含 $N$ 的子模 $L_1,L_2$。若 $x+N\in (L_1+L_2)/N$，则 $x=\ell_1+\ell_2$，其中 $\ell_i\in L_i$，故
\[
x+N=(\ell_1+N)+(\ell_2+N)\in (L_1/N)+(L_2/N).
\]
反向包含同理，故第一式成立。交的情形也直接验证：
\[
x+N\in (L_1/N)\cap(L_2/N)
\]
当且仅当 $x\in L_1$ 且 $x\in L_2$，即 $x\in L_1\cap L_2$，故第二式成立。
\end{proof}

\section{自由模与秩}

\subsection{自由模与基}

\begin{definition}
设 $R$ 为环，$M$ 为 $R$-模。若存在一族元素 $\{e_i\}_{i\in I}\subseteq M$，使得每个 $m\in M$ 都能唯一写成有限和
\[
m=\sum_{i\in I} r_ie_i\qquad (r_i\in R,\ \text{仅有限多个 }r_i\ne 0),
\]
则称 $M$ 是一个\emph{自由模}，称 $\{e_i\}_{i\in I}$ 为 $M$ 的一个\emph{基}。若 $I$ 有限且 $|I|=n$，则记
\[
M\cong R^n.
\]
\end{definition}

自由模是向量空间概念在一般环上的直接推广，但要注意：一般模未必有基，也未必能像向量空间那样分解得十分整齐。

\begin{example}
$R^n$ 配以分量加法与分量数量乘法，是自由 $R$-模，其标准基为
\[
e_1=(1,0,\dots,0),\ \dots,\ e_n=(0,\dots,0,1).
\]
任一元素 $(r_1,\dots,r_n)$ 唯一写成 $r_1e_1+\cdots+r_ne_n$。
\end{example}

\begin{example}
$\Z$-模 $\Z^n$ 是自由模，而 $\Z/n\Z$ 不是自由 $\Z$-模（当 $n\ge 2$ 时）。事实上，若 $\Z/n\Z$ 是自由 $\Z$-模，则它必与 $\Z^r$ 同构；但 $\Z^r$ 不是有限群，而 $\Z/n\Z$ 是有限群，矛盾。
\end{example}

\begin{definition}
设 $I$ 为集合。记
\[
R^{(I)}=\bigoplus_{i\in I} R.
\]
它称为以 $I$ 为指标集的\emph{自由 $R$-模}。其标准基记为 $\{e_i\}_{i\in I}$，其中 $e_i$ 在第 $i$ 个分量为 $1$，其余分量为 $0$。
\end{definition}

\subsection{自由模的泛性质}

\begin{proposition}[自由模的泛性质]
设 $F=R^{(I)}$ 是以 $\{e_i\}_{i\in I}$ 为基的自由 $R$-模，$M$ 为任意 $R$-模。若给定任意映射
\[
\phi_0\colon \{e_i\}_{i\in I}\to M,
\]
则存在唯一的 $R$-模同态
\[
\phi\colon F\to M
\]
使得对每个 $i\in I$ 都有 $\phi(e_i)=\phi_0(e_i)$。
\end{proposition}

\begin{proof}
任取 $x\in F$，可唯一写成有限和
\[
x=\sum_{i\in I} r_ie_i.
\]
定义
\[
\phi(x)=\sum_{i\in I} r_i\phi_0(e_i).
\]
由于只有有限多个 $r_i$ 非零，上式有意义；又由于表示唯一，$\phi$ 也被唯一确定。

对 $x=\sum r_ie_i$ 与 $y=\sum s_ie_i$，有
\[
\phi(x+y)=\sum (r_i+s_i)\phi_0(e_i)=\sum r_i\phi_0(e_i)+\sum s_i\phi_0(e_i)=\phi(x)+\phi(y),
\]
并且对任意 $r\in R$，
\[
\phi(rx)=\phi\Bigl(\sum (rr_i)e_i\Bigr)=\sum (rr_i)\phi_0(e_i)=r\sum r_i\phi_0(e_i)=r\phi(x).
\]
故 $\phi$ 为模同态。显然 $\phi(e_i)=\phi_0(e_i)$。唯一性来自：若 $\psi$ 也是满足条件的模同态，则
\[
\psi\Bigl(\sum r_ie_i\Bigr)=\sum r_i\psi(e_i)=\sum r_i\phi_0(e_i)=\phi\Bigl(\sum r_ie_i\Bigr),
\]
故 $\psi=\phi$。
\end{proof}

\begin{remark}
泛性质常被看作“自由模的真正定义”。它说明：从自由模到任意模的模同态，完全由基上元素的像决定。这与向量空间中“线性映射由基上的取值唯一决定”完全平行。
\end{remark}

\subsection{每个模都是自由模的商}

\begin{theorem}
任意 $R$-模 $M$ 都是某个自由 $R$-模的商模。更具体地，存在集合 $I$ 与满射模同态
\[
R^{(I)}\twoheadrightarrow M.
\]
若 $M$ 由 $n$ 个元素生成，则存在满射 $R^n\twoheadrightarrow M$。
\end{theorem}

\begin{proof}
先取集合 $I=M$ 本身。令 $F=R^{(M)}$，其基记为 $\{e_m\}_{m\in M}$。定义基上的映射
\[
e_m\longmapsto m\in M.
\]
由自由模的泛性质，存在唯一模同态
\[
\pi\colon F\to M
\]
满足 $\pi(e_m)=m$。此映射是满射，因为对任意 $m\in M$，有 $m=\pi(e_m)$。于是 $M\cong F/\Ker\pi$。

若 $M$ 由 $m_1,\dots,m_n$ 生成，则取 $F=R^n$，标准基为 $e_1,\dots,e_n$，定义
\[
\pi(e_i)=m_i.
\]
由泛性质得到模同态 $\pi\colon R^n\to M$。由生成性知 $\pi$ 满射，故 $M\cong R^n/\Ker\pi$。
\end{proof}

\subsection{秩}

在向量空间中，任意两组基有相同的基数，这个共同基数叫维数。对自由模也希望有类似概念，但对一般环需要额外说明。

\begin{theorem}[交换环上自由模的秩良定义]
设 $R$ 为非零交换环。若 $R^m\cong R^n$ 作为 $R$-模，则 $m=n$。因此有限自由 $R$-模的基的元素个数是良定义的，称为其\emph{秩}，记为 $\rank M$。
\end{theorem}

\begin{proof}
因为 $R$ 为非零交换环，存在极大理想 $\mathfrak m$。取商域样的剩余域
\[
k=R/\mathfrak m.
\]
由模同构 $R^m\cong R^n$ 可得取商后
\[
R^m/\mathfrak m R^m\cong R^n/\mathfrak m R^n.
\]
但
\[
R^m/\mathfrak m R^m\cong (R/\mathfrak m)^m=k^m,
\qquad
R^n/\mathfrak m R^n\cong (R/\mathfrak m)^n=k^n.
\]
因此 $k^m\cong k^n$ 作为 $k$-向量空间。由线性代数中维数的唯一性，得 $m=n$。
\end{proof}

\begin{remark}
若 $R$ 非交换，则“基数唯一”未必总成立；这与所谓不变基数性质（IBN）有关。对本书主要涉及的交换代数与数论背景，交换环情形已经足够。
\end{remark}

\begin{example}
理想 $(2,x)\subseteq \Z[x]$ 不是自由模的标准基那样容易理解：它作为 $\Z[x]$-模由 $2$ 与 $x$ 生成，但并非由单个元素生成，因此不是秩为 $1$ 的自由模。这个例子提醒我们：有限生成模不一定自由。
\end{example}

\section{Noether 模与 Noether 环}

\subsection{定义与等价刻画}

\begin{definition}
设 $M$ 为 $R$-模。若 $M$ 的任意递增子模链
\[
N_1\subseteq N_2\subseteq N_3\subseteq \cdots
\]
终将稳定，即存在 $N$ 使得对所有充分大的 $i$ 都有 $N_i=N$，则称 $M$ 满足\emph{升链条件}（ascending chain condition, ACC），或称 $M$ 为\emph{Noether 模}。
\end{definition}

\begin{theorem}
对 $R$-模 $M$，下列条件等价：
\begin{enumerate}[label=(\arabic*)]
  \item $M$ 满足升链条件；
  \item $M$ 的每个子模都是有限生成的；
  \item $M$ 的任意非空子模族都有极大元（按包含关系）。
\end{enumerate}
\end{theorem}

\begin{proof}
先证 $(1)\Rightarrow(2)$。设 $N\le M$。若 $N$ 不是有限生成，则可递归构造严格递增链。先任取 $x_1\in N$ 非零，设 $N_1=Rx_1$。因 $N$ 不由 $x_1$ 生成，可取 $x_2\in N\setminus N_1$，令 $N_2=Rx_1+Rx_2$。继续下去，若已构造
\[
N_k=Rx_1+\cdots+Rx_k\subsetneq N,
\]
则因 $N$ 非有限生成，可取 $x_{k+1}\in N\setminus N_k$，令
\[
N_{k+1}=N_k+Rx_{k+1}.
\]
于是得到严格递增链
\[
N_1\subsetneq N_2\subsetneq N_3\subsetneq \cdots,
\]
与 ACC 矛盾。故 $N$ 必有限生成。

再证 $(2)\Rightarrow(1)$。若存在严格递增链
\[
N_1\subsetneq N_2\subsetneq N_3\subsetneq\cdots,
\]
则其并
\[
N=\bigcup_{i\ge 1}N_i
\]
是 $M$ 的子模。由假设，$N$ 有有限生成组 $y_1,\dots,y_s$。每个 $y_j$ 落在某个 $N_{i_j}$ 中，取 $i=\max\{i_1,\dots,i_s\}$，则所有 $y_j\in N_i$，从而 $N\subseteq N_i$。这与 $N_{i+1}\subsetneq N$ 矛盾。故不存在无限严格递增链，$M$ 满足 ACC。

接着证 $(1)\Rightarrow(3)$。设 $\mathcal S$ 为非空子模族。若 $\mathcal S$ 无极大元，则任选 $N_1\in \mathcal S$。由于 $N_1$ 不是极大元，可取 $N_2\in \mathcal S$ 使 $N_1\subsetneq N_2$。继续递归下去，得到严格递增链
\[
N_1\subsetneq N_2\subsetneq N_3\subsetneq\cdots,
\]
与 ACC 矛盾。故 $\mathcal S$ 必有极大元。

最后证 $(3)\Rightarrow(1)$。设
\[
N_1\subseteq N_2\subseteq N_3\subseteq\cdots
\]
是递增链。考虑集合 $\{N_i:i\ge 1\}$，它非空，故有极大元 $N_j$。但由于链递增，对所有 $i\ge j$ 都有 $N_j\subseteq N_i$；由 $N_j$ 的极大性可知 $N_i=N_j$。故链稳定。
\end{proof}

\begin{definition}
交换环 $R$ 称为\emph{Noether 环}，若它作为 $R$-模是 Noether 模。等价地，$R$ 的每个理想都是有限生成的。
\end{definition}

\begin{example}
主理想整环都是 Noether 环，因为每个理想都由一个元素生成。特别地，$\Z$ 与 $F[x]$（$F$ 为域）都是 Noether 环。
\end{example}

\subsection{Noether 性的基本保持性质}

\begin{proposition}
若 $M$ 是 Noether 模，$N\le M$，则 $N$ 与 $M/N$ 都是 Noether 模。
\end{proposition}

\begin{proof}
先证 $N$。$N$ 的任意子模也是 $M$ 的子模，而由前述等价刻画，$M$ 的每个子模都有限生成，故 $N$ 的每个子模有限生成，从而 $N$ 是 Noether 模。

再证 $M/N$。设 $P\le M/N$。由对应定理，存在包含 $N$ 的子模 $L\le M$，使得
\[
P=L/N.
\]
因为 $L$ 是 $M$ 的子模，故有限生成，设由 $\ell_1,\dots,\ell_t$ 生成。则 $P$ 由 $\ell_1+N,\dots,\ell_t+N$ 生成。因此 $M/N$ 的每个子模有限生成，故 $M/N$ 是 Noether 模。
\end{proof}

\begin{proposition}
设
\[
0\longrightarrow M'\xrightarrow{\alpha} M\xrightarrow{\beta} M''\longrightarrow 0
\]
为正合列。若 $M'$ 与 $M''$ 都是 Noether 模，则 $M$ 也是 Noether 模。
\end{proposition}

\begin{proof}
设 $L\le M$。考虑限制映射 $\beta|_L\colon L\to M''$。其核为
\[
L\cap \Ker\beta=L\cap \Ima\alpha.
\]
由于 $\alpha$ 单射，可将 $M'$ 识别为 $\Ima\alpha$，从而 $L\cap M'$ 是 $M'$ 的子模，故有限生成。再看像 $\beta(L)$，它是 $M''$ 的子模，也有限生成。设
\[
\beta(L)=R\beta(x_1)+\cdots+R\beta(x_s),\qquad x_i\in L.
\]
又设
\[
L\cap M'=Ry_1+\cdots+Ry_t.
\]
我们断言：$L$ 由 $x_1,\dots,x_s,y_1,\dots,y_t$ 生成。任取 $x\in L$，因 $\beta(x)\in \beta(L)$，可写成
\[
\beta(x)=\sum_{i=1}^s r_i\beta(x_i)=\beta\Bigl(\sum_{i=1}^s r_ix_i\Bigr).
\]
故
\[
\beta\Bigl(x-\sum_{i=1}^s r_ix_i\Bigr)=0,
\]
于是
\[
x-\sum_{i=1}^s r_ix_i\in L\cap M'=Ry_1+\cdots+Ry_t.
\]
所以 $x$ 属于上述有限生成子模。由任意子模有限生成的刻画，$M$ 是 Noether 模。
\end{proof}

\begin{corollary}
若 $R$ 为 Noether 环，则任意有限生成 $R$-模都是 Noether 模。
\end{corollary}

\begin{proof}
设 $M$ 由 $n$ 个元素生成。由前节定理，存在满射 $R^n\twoheadrightarrow M$。因此 $M$ 是 $R^n$ 的商模。只需证明 $R^n$ 是 Noether 模。

对 $n=1$，$R^1=R$ 作为 $R$-模是 Noether 的。对一般 $n$，用归纳法。若 $R^{n-1}$ 是 Noether 模，则短正合列
\[
0\longrightarrow R^{n-1}\longrightarrow R^n\longrightarrow R\longrightarrow 0
\]
两端都 Noether，由上一命题知 $R^n$ Noether。于是其商模 $M$ 也 Noether。
\end{proof}

\subsection{Hilbert 基定理}

下面给出本章最重要的结果之一。

\begin{theorem}[Hilbert 基定理]
若 $R$ 是 Noether 环，则多项式环 $R[x]$ 也是 Noether 环。
\end{theorem}

\begin{proof}
设 $I\subseteq R[x]$ 为任意理想。我们需证明 $I$ 有限生成。

令
\[
A=\{0\}\cup\{a\in R:\exists f(x)\in I,\ f(x)\text{ 的最高次项系数为 }a\}.
\]
也就是说，$A$ 由理想 $I$ 中各多项式的首项系数组成。我们先证明 $A$ 是 $R$ 的理想。若 $a,b\in A$，则存在
\[
f(x)=ax^m+\text{低次项},\qquad g(x)=bx^n+\text{低次项}
\]
属于 $I$。于是
\[
x^n f(x)-x^m g(x)\in I.
\]
若 $a-b\ne 0$，则它的最高次项系数正是 $a-b$；若 $a-b=0$，则当然也有 $a-b\in A$。故 $a-b\in A$。对任意 $r\in R$，$rf(x)\in I$，其首项系数为 $ra$，故 $ra\in A$。因此 $A$ 为理想。

由于 $R$ Noether，理想 $A$ 有限生成。设
\[
A=(a_1,\dots,a_t).
\]
对每个 $a_i$，选择 $f_i(x)\in I$，使其首项系数为 $a_i$。设
\[
d=\max\{\deg f_i:1\le i\le t\}.
\]
接着考虑
\[
M=\{g(x)\in R[x]:\deg g<d\}\cup\{0\}.
\]
它在通常加法与数量乘法下是一个 $R$-模，并与 $R^d$ 同构（按系数向量对应）。因为 $R$ Noether，所以 $M$ 是 Noether 模。于是子模
\[
I\cap M
\]
有限生成。设
\[
I\cap M=Rg_1+\cdots+Rg_s,
\]
其中每个 $g_j$ 的次数都小于 $d$。

我们断言：
\[
I=(f_1,\dots,f_t,g_1,\dots,g_s).
\]
记右侧理想为 $J$。显然 $J\subseteq I$。下证 $I\subseteq J$。任取 $h(x)\in I$。对 $\deg h$ 作归纳。

若 $\deg h<d$，则 $h\in I\cap M$，故 $h\in J$。

设 $\deg h=n\ge d$，且已知次数小于 $n$ 的 $I$ 中元素都在 $J$ 中。设 $h$ 的首项为 $cx^n$。因为 $c\in A=(a_1,\dots,a_t)$，可写
\[
c=r_1a_1+\cdots+r_ta_t.
\]
于是多项式
\[
h_1(x)=h(x)-\sum_{i=1}^t r_ix^{n-\deg f_i}f_i(x)
\]
仍属于 $I$，且其首项已被消去，因此 $\deg h_1<n$。由归纳假设，$h_1\in J$，从而 $h\in J$。归纳完成，故 $I\subseteq J$。

因此 $I$ 有限生成。由于 $I$ 任意，$R[x]$ 是 Noether 环。
\end{proof}

\begin{corollary}
若 $R$ 为 Noether 环，则多项式环 $R[x_1,\dots,x_n]$ 对任意 $n\ge 1$ 都是 Noether 环。
\end{corollary}

\begin{proof}
对 $n$ 归纳。$n=1$ 即 Hilbert 基定理。若 $R[x_1,\dots,x_{n-1}]$ 已知是 Noether 环，则
\[
R[x_1,\dots,x_n]=R[x_1,\dots,x_{n-1}][x_n]
\]
再次应用 Hilbert 基定理即可。
\end{proof}

\begin{remark}
Hilbert 基定理是代数几何与交换代数的起点之一。它保证了多项式环中的理想、模与代数簇都可由有限数据描述。以后讨论局部化、整扩张以及模形式系数环时，它都会在背后发挥作用。
\end{remark}

\section{张量积初步}

本节默认 $R$ 为交换环，$M,N,P$ 为 $R$-模。

\subsection{平衡双线性映射与定义}

\begin{definition}
映射
\[
b\colon M\times N\to P
\]
称为\emph{$R$-双线性的}，若对每个固定的另一个变量，它都关于其中一个变量是 $R$-线性的，即对任意 $m,m'\in M$，$n,n'\in N$，$r\in R$，有
\begin{align*}
b(m+m',n)&=b(m,n)+b(m',n),\\
b(rm,n)&=r\,b(m,n),\\
b(m,n+n')&=b(m,n)+b(m,n'),\\
b(m,rn)&=r\,b(m,n).
\end{align*}
\end{definition}

\begin{definition}
$M$ 与 $N$ 的\emph{张量积}是一个 $R$-模 $M\otimes_R N$，连同一个双线性映射
\[
\tau\colon M\times N\to M\otimes_R N,
\]
满足如下泛性质：对任意 $R$-模 $P$ 与任意双线性映射 $b\colon M\times N\to P$，存在唯一的 $R$-模同态
\[
\widetilde b\colon M\otimes_R N\to P
\]
使得
\[
\widetilde b\circ \tau=b.
\]
通常记 $\tau(m,n)=m\otimes n$。
\end{definition}

这个定义与自由模的泛性质完全平行：自由模线性化“集合上的映射”，张量积则线性化“双线性映射”。

\subsection{构造}

\begin{theorem}
对任意 $R$-模 $M,N$，张量积 $M\otimes_R N$ 存在。
\end{theorem}

\begin{proof}
取集合 $M\times N$ 上的自由 $R$-模
\[
F=R^{(M\times N)}.
\]
把其标准基元记作 $[m,n]$。令 $K$ 为由下列元素生成的子模：
\begin{align*}
[m+m',n]-[m,n]-[m',n],\\
[m,n+n']-[m,n]-[m,n'],\\
[rm,n]-r[m,n],\\
[m,rn]-r[m,n],
\end{align*}
其中 $m,m'\in M$，$n,n'\in N$，$r\in R$。定义
\[
M\otimes_R N\coloneqq F/K.
\]
并令
\[
\tau\colon M\times N\to F/K,
\qquad
\tau(m,n)=[m,n]+K.
\]
记作 $m\otimes n$。

由于 $K$ 正是强制上述四类双线性关系成立的子模，故 $\tau$ 是双线性的。例如
\begin{align*}
(m+m')\otimes n
&=[m+m',n]+K\\
&=[m,n]+[m',n]+K\\
&=(m\otimes n)+(m'\otimes n),
\end{align*}
其余三条同理。

现设 $b\colon M\times N\to P$ 是任意双线性映射。先在自由模基上定义映射
\[
\phi_0([m,n])=b(m,n).
\]
由自由模的泛性质，存在唯一模同态
\[
\phi\colon F\to P
\]
使得 $\phi([m,n])=b(m,n)$。由于 $b$ 双线性，$K$ 的每个生成元都映到 $0$，故 $K\subseteq \Ker\phi$。于是 $\phi$ 可下降为唯一模同态
\[
\widetilde b\colon F/K\to P
\]
满足 $\widetilde b([m,n]+K)=b(m,n)$。即
\[
\widetilde b(m\otimes n)=b(m,n).
\]
这就证明了泛性质。

唯一性也显然：因为 $M\otimes_R N$ 由所有纯张量 $m\otimes n$ 生成，任何满足 $\widetilde b(m\otimes n)=b(m,n)$ 的模同态都只能有这一种取值。
\end{proof}

\begin{remark}
纯张量 $m\otimes n$ 的和一般不能再写成单个纯张量。也就是说，张量积中的一般元素通常是若干纯张量的有限和，而不是单个纯张量。
\end{remark}

\subsection{基本例子}

\begin{proposition}
对任意 $R$-模 $M$，有自然同构
\[
R\otimes_R M\cong M\cong M\otimes_R R.
\]
\end{proposition}

\begin{proof}
只证左边。定义双线性映射
\[
b\colon R\times M\to M,
\qquad b(r,m)=rm.
\]
由张量积的泛性质，存在唯一模同态
\[
\Phi\colon R\otimes_R M\to M,
\qquad \Phi(r\otimes m)=rm.
\]
另一方面定义
\[
\Psi\colon M\to R\otimes_R M,
\qquad \Psi(m)=1\otimes m.
\]
它是模同态，因为
\[
\Psi(rm)=1\otimes rm=r\otimes m=r(1\otimes m)=r\Psi(m).
\]
现在
\[
\Phi\Psi(m)=\Phi(1\otimes m)=m,
\]
且
\[
\Psi\Phi(r\otimes m)=\Psi(rm)=1\otimes rm=r\otimes m.
\]
故二者互逆，得到同构。
\end{proof}

\begin{proposition}
设 $I$ 为交换环 $R$ 的理想，则对任意 $R$-模 $M$，有自然同构
\[
(R/I)\otimes_R M\cong M/IM,
\]
其中
\[
IM=\Bigl\{\sum_{j=1}^t a_jm_j: a_j\in I,\ m_j\in M\Bigr\}.
\]
\end{proposition}

\begin{proof}
定义双线性映射
\[
b\colon (R/I)\times M\to M/IM,
\qquad b(r+I,m)=rm+IM.
\]
需证良定义：若 $r-r'\in I$，则
\[
rm-r'm=(r-r')m\in IM,
\]
故 $rm+IM=r'm+IM$。于是由泛性质得到模同态
\[
\Phi\colon (R/I)\otimes_R M\to M/IM,
\qquad \Phi((r+I)\otimes m)=rm+IM.
\]
再定义
\[
\Psi\colon M\to (R/I)\otimes_R M,
\qquad \Psi(m)=(1+I)\otimes m.
\]
若 $x=\sum a_jm_j\in IM$，则
\[
\Psi(x)=\sum (1+I)\otimes a_jm_j=\sum (a_j+I)\otimes m_j=0,
\]
因为 $a_j\in I$。故 $IM\subseteq \Ker\Psi$，于是 $\Psi$ 下降为模同态
\[
\overline\Psi\colon M/IM\to (R/I)\otimes_R M,
\qquad \overline\Psi(m+IM)=(1+I)\otimes m.
\]
现在直接计算：
\[
\Phi\overline\Psi(m+IM)=\Phi((1+I)\otimes m)=m+IM,
\]
而对纯张量有
\[
\overline\Psi\Phi((r+I)\otimes m)=\overline\Psi(rm+IM)=(1+I)\otimes rm=(r+I)\otimes m.
\]
故两映射互逆，得所求同构。
\end{proof}

\begin{corollary}
对任意阿贝尔群 $A$ 与正整数 $n$，有
\[
(\Z/n\Z)\otimes_\Z A\cong A/nA.
\]
特别地，
\[
\Q\otimes_\Z (\Z/n\Z)=0.
\]
\end{corollary}

\begin{proof}
第一式是将上命题取 $R=\Z$、$I=n\Z$ 的特例。第二式可由第一式取 $A=\Q$ 得到，因为 $n\Q=\Q$，故
\[
\Q/n\Q=0.
\]
\end{proof}

\begin{example}
若 $V,W$ 是域 $F$ 上有限维向量空间，且 $\dim_FV=m$、$\dim_FW=n$，则
\[
V\otimes_F W\cong F^{mn}.
\]
事实上，取基 $v_1,\dots,v_m$ 与 $w_1,\dots,w_n$，则 $v_i\otimes w_j$ 生成张量积，并且可证它们线性无关，因此构成一组基。
\end{example}

\subsection{函子性与基本性质}

\begin{proposition}[函子性]
设 $f\colon M\to M'$、$g\colon N\to N'$ 为 $R$-模同态，则存在唯一模同态
\[
f\otimes g\colon M\otimes_R N\to M'\otimes_R N'
\]
满足
\[
(f\otimes g)(m\otimes n)=f(m)\otimes g(n).
\]
\end{proposition}

\begin{proof}
映射
\[
M\times N\to M'\otimes_R N',
\qquad (m,n)\mapsto f(m)\otimes g(n)
\]
是双线性的，因为 $f,g$ 都线性。由张量积的泛性质，存在唯一模同态 $f\otimes g$ 使上式成立。
\end{proof}

\begin{proposition}
对任意模族 $\{M_i\}_{i\in I}$ 与模 $N$，有自然同构
\[
\Bigl(\bigoplus_{i\in I} M_i\Bigr)\otimes_R N\cong \bigoplus_{i\in I}(M_i\otimes_R N).
\]
\end{proposition}

\begin{proof}
记 $M=\bigoplus_i M_i$。对 $(m_i)_i\in M$ 与 $n\in N$，定义
\[
\theta((m_i)_i,n)=(m_i\otimes n)_i\in \bigoplus_i(M_i\otimes_R N).
\]
因 $(m_i)_i$ 只有有限多个非零分量，故右侧确实落在直和中。且 $\theta$ 显然双线性。由泛性质得模同态
\[
\widetilde\theta\colon M\otimes_R N\to \bigoplus_i(M_i\otimes_R N).
\]
另一方面，每个嵌入 $\iota_i\colon M_i\hookrightarrow M$ 诱导
\[
\iota_i\otimes \Id_N\colon M_i\otimes_R N\to M\otimes_R N.
\]
由直和的泛性质，这些映射合并成唯一模同态
\[
\eta\colon \bigoplus_i(M_i\otimes_R N)\to M\otimes_R N.
\]
在纯张量上检查可知二者互逆：
\[
\eta\widetilde\theta((m_i)_i\otimes n)=\sum_i \iota_i(m_i)\otimes n=(m_i)_i\otimes n,
\]
而对来自第 $i$ 个分量的纯张量 $m_i\otimes n$，有
\[
\widetilde\theta\eta(m_i\otimes n)=m_i\otimes n.
\]
故得到所求同构。
\end{proof}

\begin{theorem}[右正合性]
设
\[
L\xrightarrow{u} M\xrightarrow{v} N\to 0
\]
是正合列，则对任意 $R$-模 $P$，列
\[
L\otimes_R P\xrightarrow{u\otimes \Id_P} M\otimes_R P\xrightarrow{v\otimes \Id_P} N\otimes_R P\to 0
\]
也是正合的。
\end{theorem}

\begin{proof}
先证末端满射。因为 $v$ 满射，对任意纯张量 $n\otimes p\in N\otimes_R P$，可取 $m\in M$ 使 $v(m)=n$，则
\[
(v\otimes \Id_P)(m\otimes p)=v(m)\otimes p=n\otimes p.
\]
而 $N\otimes_R P$ 由纯张量生成，故 $v\otimes \Id_P$ 满射。

再证
\[
\Ima(u\otimes \Id_P)\subseteq \Ker(v\otimes \Id_P).
\]
这是因为
\[
(v\otimes \Id_P)\circ (u\otimes \Id_P)=(v\circ u)\otimes \Id_P=0.
\]

关键是证明反向包含。设 $Q=(M\otimes_R P)/\Ima(u\otimes \Id_P)$，记商映射为 $q\colon M\otimes_R P\to Q$。定义双线性映射
\[
b\colon M\times P\to Q,
\qquad b(m,p)=q(m\otimes p).
\]
若 $m-m'=u(\ell)$ 属于 $\Ima u$，则
\[
q(m\otimes p)-q(m'\otimes p)=q((m-m')\otimes p)=q((u\otimes \Id_P)(\ell\otimes p))=0,
\]
故 $b(m,p)$ 只依赖于 $v(m)\in N$。因此可以定义
\[
\overline b\colon N\times P\to Q,
\qquad \overline b(v(m),p)=q(m\otimes p).
\]
它良定义且双线性。由张量积的泛性质，存在唯一模同态
\[
\gamma\colon N\otimes_R P\to Q
\]
满足
\[
\gamma(v(m)\otimes p)=q(m\otimes p).
\]
于是
\[
\gamma\circ (v\otimes \Id_P)=q.
\]
若 $x\in \Ker(v\otimes \Id_P)$，则
\[
q(x)=\gamma((v\otimes \Id_P)(x))=0.
\]
故 $x\in \Ker q=\Ima(u\otimes \Id_P)$。因此
\[
\Ker(v\otimes \Id_P)=\Ima(u\otimes \Id_P).
\]
正合性得证。
\end{proof}

\begin{remark}
张量积一般\emph{不左正合}。例如正合列
\[
0\to \Z\xrightarrow{\times 2} \Z\to \Z/2\Z\to 0
\]
张量 $\Z/2\Z$ 后得到
\[
0\to \Z/2\Z\xrightarrow{0} \Z/2\Z\to \Z/2\Z\to 0,
\]
中间箭头不再单射。正因为如此，平坦模的概念才显得重要。
\end{remark}

\subsection{扩张标量}

\begin{definition}
设 $R\to S$ 为交换环同态，$M$ 为 $R$-模。则
\[
S\otimes_R M
\]
称为把 $M$ 从 $R$ 扩张到 $S$ 后得到的模，或称为\emph{扩张标量}。其 $S$-模结构由
\[
s'\cdot (s\otimes m)=(s's)\otimes m
\]
给出。
\end{definition}

\begin{proposition}
上述定义确实给出 $S\otimes_R M$ 的一个 $S$-模结构；并且对任意 $S$-模 $N$，有自然同构
\[
\Hom_S(S\otimes_R M,N)\cong \Hom_R(M,N),
\]
其中右边把 $N$ 视为 $R$-模。
\end{proposition}

\begin{proof}
$S$-模公理在纯张量上直接验证。例如
\[
(s_1s_2)\cdot (s\otimes m)=(s_1s_2s)\otimes m=s_1\cdot ((s_2s)\otimes m)=s_1\cdot (s_2\cdot (s\otimes m)).
\]
其余公理同理。

接着构造同构。若 $\Phi\colon S\otimes_R M\to N$ 为 $S$-模同态，则定义
\[
\alpha(\Phi)(m)=\Phi(1\otimes m).
\]
这是 $R$-模同态，因为
\[
\alpha(\Phi)(rm)=\Phi(1\otimes rm)=\Phi(r\otimes m)=r\Phi(1\otimes m)=r\alpha(\Phi)(m).
\]
反过来，若 $f\colon M\to N$ 为 $R$-模同态，定义双线性映射
\[
b\colon S\times M\to N,
\qquad b(s,m)=s f(m),
\]
这里右侧用的是 $N$ 的 $S$-模结构。由泛性质得到唯一 $R$-模同态
\[
\widetilde f\colon S\otimes_R M\to N,
\qquad \widetilde f(s\otimes m)=sf(m).
\]
并且它实际上是 $S$-线性的，因为
\[
\widetilde f(s'\cdot (s\otimes m))=\widetilde f((s's)\otimes m)=s'sf(m)=s'\widetilde f(s\otimes m).
\]
故得到映射
\[
\beta\colon \Hom_R(M,N)\to \Hom_S(S\otimes_R M,N),\qquad \beta(f)=\widetilde f.
\]
最后验证互逆：
\[
\alpha(\beta(f))(m)=\widetilde f(1\otimes m)=f(m),
\]
而对 $\Phi\in \Hom_S(S\otimes_R M,N)$ 与纯张量 $s\otimes m$，有
\[
\beta(\alpha(\Phi))(s\otimes m)=s\alpha(\Phi)(m)=s\Phi(1\otimes m)=\Phi(s\otimes m).
\]
故 $\alpha$ 与 $\beta$ 互逆。
\end{proof}

\begin{example}
若 $M$ 是阿贝尔群，则
\[
\Q\otimes_\Z M
\]
可看作“把 $M$ 中的整数倍关系有理化”。若 $M$ 有扭部分 $T$，则 $\Q\otimes_\Z T=0$；若 $M\cong \Z^r\oplus T$，则
\[
\Q\otimes_\Z M\cong \Q^r.
\]
因此 $\Q\otimes_\Z M$ 只保留自由部分的秩信息。
\end{example}

\section{模形式空间作为 Hecke 模}

本节说明模论如何进入模形式理论。我们不系统发展 Hecke 算子的解析定义，而把重点放在“模结构”这个代数视角上。

\subsection{Hecke 代数的作用}

设 $\Gamma$ 是 $\SL_2(\Z)$ 的一个同余子群，$k\ge 0$ 为整数。模形式空间 $\Mk_k(\Gamma)$ 与尖点模形式空间 $\Sk_k(\Gamma)$ 都是有限维 $\C$-向量空间。对每个正整数 $n$，存在一个 $\C$-线性算子
\[
T_n\colon \Mk_k(\Gamma)\to \Mk_k(\Gamma),
\]
称为第 $n$ 个 Hecke 算子；并且 $T_n$ 保持 $\Sk_k(\Gamma)$。

\begin{definition}
由所有 Hecke 算子生成的 $\C$-子代数
\[
\mathbb T_k(\Gamma)\subseteq \End_\C(\Mk_k(\Gamma))
\]
称为权 $k$、群 $\Gamma$ 的\emph{Hecke 代数}。若只作用在尖点模形式上，也得到
\[
\mathbb T_k^{\mathrm{cusp}}(\Gamma)\subseteq \End_\C(\Sk_k(\Gamma)).
\]
\end{definition}

\begin{proposition}
$\Mk_k(\Gamma)$ 是一个 $\mathbb T_k(\Gamma)$-模，$\Sk_k(\Gamma)$ 是一个 $\mathbb T_k^{\mathrm{cusp}}(\Gamma)$-模。
\end{proposition}

\begin{proof}
把 $\mathbb T_k(\Gamma)$ 看作作用于 $\Mk_k(\Gamma)$ 的算子环的子环。对 $A\in \mathbb T_k(\Gamma)$ 与 $f\in \Mk_k(\Gamma)$，定义
\[
A\cdot f\coloneqq A(f).
\]
由于 $\mathbb T_k(\Gamma)$ 是 $\End_\C(\Mk_k(\Gamma))$ 的子环，算子加法与合成满足
\begin{align*}
(A+B)\cdot f&=A(f)+B(f)=A\cdot f+B\cdot f,\\
A\cdot (f+g)&=A(f+g)=A(f)+A(g),\\
(AB)\cdot f&=A(B(f))=A\cdot (B\cdot f),\\
1\cdot f&=f.
\end{align*}
故模公理成立，因此 $\Mk_k(\Gamma)$ 成为 $\mathbb T_k(\Gamma)$-模。对 $\Sk_k(\Gamma)$ 的证明完全相同。
\end{proof}

\begin{remark}
这个命题本身形式上很简单，但意义重大：它把模形式空间从一个“函数空间”变成了一个“代数表示对象”。此后关于本征形式、特征值、同余与 Galois 表示的许多问题，都可理解为研究这个 Hecke 模的结构。
\end{remark}

\subsection{$q$-展开与 Hecke 算子}

对全模群 $\Gamma=\SL_2(\Z)$，若
\[
f(q)=\sum_{n=0}^\infty a_n q^n\in \Mk_k(\SL_2(\Z)),
\]
则对素数 $p$，Hecke 算子 $T_p$ 在 $q$-展开上的作用为
\[
T_p f=\sum_{n=0}^\infty \bigl(a_{pn}+p^{k-1}a_{n/p}\bigr)q^n,
\]
其中约定当 $p\nmid n$ 时 $a_{n/p}=0$。

\begin{example}
设 $E_k$ 是标准归一化 Eisenstein 级数（$k\ge 4$ 且为偶数），则
\[
T_pE_k=(1+p^{k-1})E_k.
\]
这是因为 $E_k$ 的傅里叶系数满足经典公式
\[
a_n(E_k)=\sigma_{k-1}(n)=\sum_{d\mid n} d^{k-1},
\]
而 Hecke 算子的作用给出递推
\[
a_n(T_pf)=a_{pn}(f)+p^{k-1}a_{n/p}(f).
\]
对 $f=E_k$，利用约数和函数的乘法性质，便得到上式。因此 $E_k$ 是 $T_p$ 的本征向量，本征值为 $1+p^{k-1}$。
\end{example}

\subsection{固定一个 Hecke 算子时的模论解释}

Hecke 代数通常由许多彼此可交换的算子组成。若先固定其中一个算子 $T\in \End_\C(\Mk_k(\Gamma))$，则由 \S10.1 中 $\C[x]$-模与线性算子的对应，$\Mk_k(\Gamma)$ 可以视为一个有限生成 $\C[x]$-模，其中 $x$ 的作用定义为
\[
x\cdot f=T(f).
\]

由于 $\Mk_k(\Gamma)$ 是有限维 $\C$-向量空间，这个 $\C[x]$-模是有限生成且扭的。于是可以应用 PID 上有限生成模的结构定理：它分解为若干循环模的直和，亦即
\[
\Mk_k(\Gamma)\cong \bigoplus_{i=1}^r \C[x]/\bigl((x-\lambda_i)^{e_i}\bigr).
\]
这正是 Jordan 标准形在模论语言中的表达。特别地：
\begin{itemize}
  \item 若所有 $e_i=1$，则 $T$ 可对角化；
  \item 若某个 $e_i>1$，则出现非平凡 Jordan 块。
\end{itemize}

因此，对 Hecke 算子的谱分解可以理解为对相应 $\C[x]$-模的结构分析。

\subsection{重数一现象}

在很多重要情形下，Hecke 模的结构比一般线性代数更刚性。一个典型结果是重数一（multiplicity one）定理。

\begin{theorem}[重数一定理，只陈述]
设 $f,g\in \Sk_k(\SL_2(\Z))$ 是两个归一化 Hecke 本征尖点模形式。若对每个素数 $p$ 都有
\[
T_pf=a_pf,\qquad T_pg=a_pg,
\]
并且二者的本征值相同，则 $f=g$。

更一般地，只要两个归一化本征形式具有相同的 Hecke 本征值系统，它们就必须相等。
\end{theorem}

\begin{remark}
从模论观点看，重数一定理说明某些 Hecke 模的本征空间不会“以不可控的重数重复出现”。这正是把傅里叶系数、L-函数与 Galois 表示联系起来的关键刚性来源之一。严格证明需要更深入的模形式理论，超出本章范围。
\end{remark}

\section*{习题}

下面共 $55$ 题。基础题用 \basic{} 标记，进阶题用 \medium{} 标记，挑战题用 \hard{} 标记。

\subsection*{基础题}

\begin{enumerate}[label=\textbf{10.\arabic*.}, leftmargin=3.2em]
  \item (\basic) 设 $A$ 为阿贝尔群，验证它按整数倍运算成为一个 $\Z$-模。

  \item (\basic) 设 $F$ 为域，$V$ 为 $F$-向量空间。逐条检查：$V$ 的向量空间公理正好给出 $F$-模公理。

  \item (\basic) 设 $R$ 为交换环，$I\subseteq R$ 为理想。证明 $I$ 是 $R$-模。

  \item (\basic) 设 $M,N$ 为 $R$-模。证明直积 $M\times N$ 及直和 $M\oplus N$ 都是 $R$-模。

  \item (\basic) 设 $f\colon M\to N$ 为模同态。证明 $f(0)=0$，并证明 $f(-m)=-f(m)$。

  \item (\basic) 求出 $\Z$ 作为 $\Z$-模的全部子模。

  \item (\basic) 求出 $\Z/12\Z$ 作为 $\Z$-模的全部子模。

  \item (\basic) 设 $R$ 为交换环，$I,J$ 为理想。证明 $I\cap J$ 与 $I+J$ 都是 $R$ 的子模。

  \item (\basic) 在 $\Z$-模 $\Z^2$ 中，证明集合
  \[
  N=\{(a,b)\in \Z^2: a\equiv b\pmod 2\}
  \]
  是子模。

  \item (\basic) 设 $N=n\Z\le \Z$。写出商模 $\Z/N$ 的全部元素，并计算其加法表（取 $n=4$）。

  \item (\basic) 在 $\Z^2$ 中取子模 $N=\langle (1,1)\rangle$。证明 $\Z^2/N\cong \Z$。

  \item (\basic) 设 $M=R$，$N=I$ 为理想。说明商模 $M/N$ 与商环 $R/I$ 作为 $R$-模一致。

  \item (\basic) 设 $f\colon \Z^2\to \Z$，$f(a,b)=a-b$。求 $\Ker f$ 与 $\Ima f$，并验证第一同构定理。

  \item (\basic) 设 $R$ 为环，$A,B\le M$。证明 $A+B$ 是包含 $A$ 与 $B$ 的最小子模。

  \item (\basic) 设 $R$ 为交换环。证明 $R^n$ 是自由模，并写出其标准基。

  \item (\basic) 证明 $\Z/n\Z$（$n\ge 2$）不是自由 $\Z$-模。

  \item (\basic) 证明任意循环 $R$-模都同构于某个商模 $R/I$。

  \item (\basic) 设 $M$ 由 $m_1,m_2,m_3$ 生成。构造一个满射 $R^3\twoheadrightarrow M$。

  \item (\basic) 设 $F$ 为域，证明有限维 $F$-向量空间都是有限秩自由 $F$-模。

  \item (\basic) 设 $R$ 为 Noether 环，$M$ 为有限生成 $R$-模。解释为什么 $M$ 必是 Noether 模。

  \item (\basic) 证明若 $M$ 是 Noether 模，则任意商模 $M/N$ 也是 Noether 模。

  \item (\basic) 计算 $\Z\otimes_\Z M$，其中 $M$ 为任意阿贝尔群。

  \item (\basic) 计算 $(\Z/n\Z)\otimes_\Z \Z/m\Z$。

  \item (\basic) 设 $V,W$ 为域 $F$ 上向量空间，且基分别为 $\{v_i\}$ 与 $\{w_j\}$。证明 $\{v_i\otimes w_j\}$ 生成 $V\otimes_F W$。

  \item (\basic) 设 $f(q)=\sum_{n\ge 0} a_nq^n\in \Mk_k(\SL_2(\Z))$。根据 Hecke 算子公式，写出 $T_2f$ 的前五个傅里叶系数。
\end{enumerate}

\subsection*{进阶题}

\begin{enumerate}[label=\textbf{10.\arabic*.}, leftmargin=3.2em, start=26]
  \item (\medium) 设 $f\colon M\to N$ 为模同态。给出第一同构定理的另一种证明：直接构造 $M/\Ker f$ 到 $\Ima f$ 的双射并验证其线性性。

  \item (\medium) 证明第二同构定理与第三同构定理。

  \item (\medium) 证明对应定理，并说明它如何推广群论中的正规子群对应定理。

  \item (\medium) 设 $M$ 为 $R$-模，$A,B,C\le M$。证明
  \[
  A\cap (B+C)\supseteq (A\cap B)+(A\cap C),
  \]
  并给出不取等号的例子。

  \item (\medium) 设 $R$ 为交换环，证明：若 $R^m\cong R^n$，则对任意极大理想 $\mathfrak m$ 都有 $(R/\mathfrak m)^m\cong (R/\mathfrak m)^n$，从而 $m=n$。

  \item (\medium) 设 $I=(x,y)\subseteq k[x,y]$。证明 $I$ 不是自由模 $k[x,y]$ 的秩 $1$ 子模。

  \item (\medium) 证明：若 $M$ 与 $N$ 都是 Noether 模，则 $M\oplus N$ 也是 Noether 模。

  \item (\medium) 证明：交换环 $R$ 是 Noether 环，当且仅当 $R$ 的每个理想都有限生成。

  \item (\medium) 设 $R$ 为 Noether 环。利用 Hilbert 基定理证明 $R[x_1,\dots,x_n]$ 是 Noether 环。

  \item (\medium) 证明 $\Z[x]$ 是 Noether 环，但不是主理想整环。

  \item (\medium) 设 $R$ 为 PID，$M$ 为有限生成扭 $R$-模。证明若
  \[
  M\cong R/(a)\oplus R/(b),\qquad (a,b)=1,
  \]
  则 $M\cong R/(ab)$。

  \item (\medium) 利用有限生成阿贝尔群结构定理，分类所有阶为 $72$ 的阿贝尔群。

  \item (\medium) 设 $A$ 为阿贝尔群。证明
  \[
  (\Z/n\Z)\otimes_\Z A\cong A/nA.
  \]
  并据此计算 $(\Z/6\Z)\otimes_\Z \Z/15\Z$。

  \item (\medium) 证明张量积满足交换性：对交换环 $R$，构造自然同构
  \[
  M\otimes_R N\cong N\otimes_R M.
  \]

  \item (\medium) 证明张量积满足结合律：构造自然同构
  \[
  (M\otimes_R N)\otimes_R P\cong M\otimes_R (N\otimes_R P).
  \]

  \item (\medium) 设
  \[
  L\to M\to N\to 0
  \]
  正合。用张量积的泛性质重新证明右正合性。

  \item (\medium) 给出一个例子说明张量积不保持单射。

  \item (\medium) 设 $S$ 是交换环 $R$ 的一个代数。证明扩张标量函子 $M\mapsto S\otimes_R M$ 保持有限直和。

  \item (\medium) 设 $T\in \End_F(V)$，其中 $V$ 为有限维 $F$-向量空间。利用 $F[x]$-模观点说明最小多项式为何湮没整个模。

  \item (\medium) 设 $f\in \Sk_k(\SL_2(\Z))$ 是归一化 Hecke 本征形式，且
  \[
  f(q)=\sum_{n\ge 1} a_nq^n.
  \]
  证明若 $T_pf=a_pf$，则对任意素数 $p$，有
  \[
  a_p(f)=a_p.
  \]
\end{enumerate}

\subsection*{挑战题}

\begin{enumerate}[label=\textbf{10.\arabic*.}, leftmargin=3.2em, start=46]
  \item (\hard) 设 $R$ 为局部环，极大理想为 $\mathfrak m$，$M$ 为有限生成 $R$-模。证明 Nakayama 引理：若 $\mathfrak m M=M$，则 $M=0$。

  \item (\hard) 在上一题条件下，证明 Nakayama 引理的推论：若 $x_1,\dots,x_n$ 在 $M/\mathfrak m M$ 中的像生成该商模，则 $x_1,\dots,x_n$ 生成 $M$。

  \item (\hard) 定义：若对每个短正合列
  \[
  0\to A\to B\to C\to 0
  \]
  张量 $M$ 后仍保持正合，则称 $M$ 为平坦模。证明任意自由模都是平坦模。

  \item (\hard) 设 $R$ 为整环。定义 $R$-模 $M$ 的扭子模
  \[
  M_{\mathrm{tor}}=\{m\in M:\exists 0\ne r\in R,\ rm=0\}.
  \]
  证明 $M_{\mathrm{tor}}$ 是 $M$ 的子模。

  \item (\hard) 设 $R$ 为 PID，$M$ 为有限生成 $R$-模。证明 $M\cong R^r\oplus M_{\mathrm{tor}}$，并说明 $r$ 与 $M_{\mathrm{tor}}$ 唯一确定。

  \item (\hard) 设 $R$ 为交换环，$I\subseteq \Rad(R)$，$M$ 为有限生成 $R$-模。证明若 $M=IM$，则 $M=0$。将其视为 Nakayama 引理的一种形式。

  \item (\hard) 设 $V$ 为有限维 $\C$-向量空间，$T\in \End_\C(V)$。把 $V$ 视为 $\C[x]$-模，证明：$V$ 是扭模，且其湮没子理想由 $T$ 的最小多项式生成。

  \item (\hard) 设 $\Gamma=\SL_2(\Z)$。验证 Hecke 代数 $\mathbb T_k(\Gamma)$ 对 $\Mk_k(\Gamma)$ 的作用满足模公理。

  \item (\hard) 设
  \[
  f(q)=\sum_{n\ge 0} a_nq^n\in \Mk_k(\SL_2(\Z)).
  \]
  用 Hecke 算子公式证明：若 $f$ 是 $T_p$ 的本征向量，本征值为 $\lambda_p$，则
  \[
  a_p=\lambda_p a_1.
  \]
  特别地，当 $f$ 归一化（即 $a_1=1$）时，$a_p=\lambda_p$。

  \item (\hard) 设 $V=\Mk_k(\Gamma)$，固定一个 Hecke 算子 $T$。把 $V$ 视为 $\C[x]$-模。说明为什么 $T$ 可对角化当且仅当该 $\C[x]$-模可分解为若干简单模 $\C[x]/(x-\lambda)$ 的直和。
\end{enumerate}
